\documentclass{article}
\usepackage{amsmath,amssymb,amsthm}
\usepackage{algorithm}
\usepackage{algpseudocode}
\usepackage{enumitem}
\newtheorem{theorem}{Theorem}
\newtheorem{lemma}{Lemma}
\newtheorem{assumption}{Assumption}
\newcommand{\E}{\mathbb{E}}
\newcommand{\R}{\mathbb{R}}
\newcommand{\norm}[1]{\left\|#1\right\|}
\newcommand{\topk}[1]{\mathcal{T}_K\!\left(#1\right)}
\begin{document}

\section*{Top‑K Stochastic Gradient Descent (Top‑K SGD)}

\section*{Mathematical Formulation}
Let $f:\R^{d}\to\R$ be the (possibly non‑convex) objective and let $z$ denote a random data sample.
At iteration $t$ a stochastic gradient $g_t$ is computed
\[
g_t \;:=\; \nabla f(w_t;z_t)\;\in\R^{d}.
\]
Define the \emph{top‑$K$ operator}
\[
\topk{g}\;:=\;\operatorname{argmin}_{\tilde g\in\R^{d}}
\bigl\{ \norm{g-\tilde g}^2 \;:\; \tilde g_i =0\;\text{for all but at most $K$ indices}\bigr\},
\]
i.e. $\topk{g}$ keeps the $K$ components of $g$ with largest absolute value and zeros out the rest.
The update rule of Top‑K SGD is
\begin{equation}
\label{eq:update}
w_{t+1}= w_t - \eta_t \,\topk{g_t},
\qquad t=0,1,\dots,T-1,
\end{equation}
where $\eta_t>0$ is a stepsize (learning rate).  

\section*{Pseudocode}
\begin{algorithm}[H]
\caption{Top‑K SGD}
\begin{algorithmic}[1]
\State \textbf{Input:} initial point $w_0\in\R^{d}$, stepsize schedule $\{\eta_t\}_{t\ge0}$, sparsity level $K\le d$.
\For{$t=0$ \textbf{to} $T-1$}
    \State Sample data $z_t$ and compute stochastic gradient $g_t=\nabla f(w_t;z_t)$.
    \State $\displaystyle \tilde g_t \gets \topk{g_t}$ \Comment{keep $K$ largest magnitudes}
    \State $w_{t+1} \gets w_t - \eta_t \tilde g_t$.
\EndFor
\State \textbf{Output:} $\{w_t\}_{t=0}^{T}$ (or the averaged iterate $\bar w_T=\frac1T\sum_{t=0}^{T-1}w_t$).
\end{algorithmic}
\end{algorithm}

\section*{Dry Run Example}
Consider the scalar quadratic $f(w)=(w-3)^2$, thus $d=1$ and the true gradient is $\nabla f(w)=2(w-3)$.
Take $w_0=0$, stepsize $\eta=0.1$, and $K=1$ (the only coordinate is always kept).

\begin{center}
\begin{tabular}{c|c|c|c}
$t$ & $g_t=\nabla f(w_t)$ & $w_{t+1}=w_t-\eta g_t$ & $f(w_{t+1})$\\\hline
0 & $2(0-3)=-6$ & $0-0.1(-6)=0.6$ & $(0.6-3)^2=5.76$\\
1 & $2(0.6-3)=-4.8$ & $0.6-0.1(-4.8)=1.08$ & $(1.08-3)^2=3.6864$\\
2 & $2(1.08-3)=-3.84$ & $1.08-0.1(-3.84)=1.464$ & $(1.464-3)^2=2.3660$
\end{tabular}
\end{center}
The iterates move toward the minimiser $w^\star=3$ and the objective values decrease.

\newpage
\section*{Assumptions}
\begin{assumption}[Smoothness]
\label{as:smooth}
$f$ is $L$‑smooth: for all $x,y\in\R^{d}$,
\[
f(y)\le f(x)+\langle \nabla f(x),y-x\rangle+\frac{L}{2}\norm{y-x}^{2}.
\]
\end{assumption}
\begin{assumption}[Unbiased Stochastic Gradient and Bounded Variance]
\label{as:unbiased}
For every $t$, conditioned on $w_t$,
\[
\E[g_t\mid w_t]=\nabla f(w_t),\qquad 
\E\bigl[\norm{g_t-\nabla f(w_t)}^{2}\mid w_t\bigr]\le\sigma^{2}.
\]
\end{assumption}
\begin{assumption}[Top‑K Compression Error]
\label{as:topk}
For any vector $u\in\R^{d}$,
\[
\E\!\left[\norm{\topk{u}-u}^{2}\right]\le\left(1-\frac{K}{d}\right)\norm{u}^{2}.
\]
(The expectation is taken over the (possibly random) tie‑breaking in $\topk{\cdot}$.)
\end{assumption}
\begin{assumption}[Stepsize]
\label{as:stepsize}
The stepsize is constant $\eta_t\equiv\eta$ with
\[
0<\eta\le \frac{1}{L\sqrt{1-\frac{K}{d}}}.
\]
\end{assumption}

\section*{Main Convergence Theorem}
\begin{theorem}[Convergence of Top‑K SGD for Non‑Convex $L$‑smooth Objectives]
\label{thm:main}
Under Assumptions \ref{as:smooth}–\ref{as:stepsize} and after $T$ iterations,
\[
\frac{1}{T}\sum_{t=0}^{T-1}\E\!\bigl[\norm{\nabla f(w_t)}^{2}\bigr]
\;\le\;
\frac{2\bigl(f(w_0)-f^\star\bigr)}{\eta T}
\;+\;
\eta L\sigma^{2}
\;+\;
\frac{L\eta\sigma^{2}}{K/d},
\]
where $f^\star=\inf_{w}f(w)$.  In particular, choosing $\eta=c/\sqrt{T}$ for a suitable constant $c$ yields
\[
\frac{1}{T}\sum_{t=0}^{T-1}\E\!\bigl[\norm{\nabla f(w_t)}^{2}\bigr]
= O\!\left(\frac{1}{\sqrt{T}}\right).
\]
\end{theorem}

\section*{Theoretical Properties}
\begin{itemize}[leftmargin=*]
\item \textbf{Stability.} The bound requires $\eta$ to be at most $1/(L\sqrt{1-K/d})$, guaranteeing that the compression error does not destabilise the iteration.
\item \textbf{Hyperparameter Sensitivity.} Larger $K$ (less sparsification) reduces the factor $1-\frac{K}{d}$, tightening the bound; conversely, very small $K$ slows convergence through the additional $\frac{L\eta\sigma^{2}}{K/d}$ term.
\item \textbf{Comparison to Full‑SGD.} Setting $K=d$ makes $\topk{g}=g$, the compression term vanishes and Theorem~\ref{thm:main} reduces to the classic $O(1/\sqrt{T})$ result for SGD.
\end{itemize}

\section*{Discussion}
The theorem shows that even with aggressive sparsification (small $K$) the algorithm retains the same $O(1/\sqrt{T})$ rate as full SGD up to a multiplicative constant that depends on the compression ratio $K/d$. This matches empirical observations that Top‑K communication substantially reduces bandwidth while preserving convergence speed for deep learning workloads.

\newpage
\section*{Preliminaries (Definitions and Lemmas)}
\begin{lemma}[Smoothness Inequality]
\label{lem:smooth}
If $f$ is $L$‑smooth, then for any $x,y\in\R^{d}$,
\[
f(y)\le f(x)+\langle \nabla f(x),y-x\rangle+\frac{L}{2}\norm{y-x}^{2}.
\]
\end{lemma}
\begin{proof}
Standard consequence of the definition of $L$‑smoothness; see e.g. \cite{nesterov2004introductory}. \hfill\qedsymbol
\end{proof}

\begin{lemma}[Compression Error of Top‑K]
\label{lem:topk}
For any $u\in\R^{d}$,
\[
\E\!\left[\norm{\topk{u}-u}^{2}\right]\le\left(1-\frac{K}{d}\right)\norm{u}^{2}.
\]
\end{lemma}
\begin{proof}
Let $S\subseteq\{1,\dots,d\}$ be the index set of the $K$ largest absolute entries of $u$. Then
\[
\topk{u}_i=
\begin{cases}
u_i,& i\in S,\\
0,& i\notin S.
\end{cases}
\]
Hence
\[
\norm{\topk{u}-u}^{2}
= \sum_{i\notin S}u_i^{2}
\le \frac{d-K}{d}\sum_{i=1}^{d}u_i^{2}
= \Bigl(1-\frac{K}{d}\Bigr)\norm{u}^{2}.
\]
Taking expectation over any random tie‑breaking does not increase the RHS. \hfill\qedsymbol
\end{proof}

\section*{Main Proof (Step‑by‑Step Derivation)}
We prove Theorem~\ref{thm:main}. All expectations are taken with respect to the randomness of the stochastic gradients and the possible randomness of the top‑$K$ operator.

\noindent\textbf{Notation.} Define the compression error
\[
e_t := \topk{g_t}-g_t,\qquad \tilde g_t:=\topk{g_t}=g_t+e_t .
\]

\noindent\textbf{Step 1.} Apply Lemma~\ref{lem:smooth} with $x=w_t$, $y=w_{t+1}=w_t-\eta\tilde g_t$.
\[
\begin{aligned}
f(w_{t+1})
&\le f(w_t)+\langle \nabla f(w_t),\, -\eta\tilde g_t\rangle
   +\frac{L}{2}\eta^{2}\norm{\tilde g_t}^{2}. \tag{1}
\end{aligned}
\]
[Step 1]

\noindent\textbf{Step 2.} Decompose $\tilde g_t$ into $g_t+e_t$ and use unbiasedness $\E[g_t\mid w_t]=\nabla f(w_t)$.
\[
\begin{aligned}
\langle \nabla f(w_t),\tilde g_t\rangle
&= \langle \nabla f(w_t),g_t\rangle+\langle \nabla f(w_t),e_t\rangle.
\end{aligned}
\]
Take conditional expectation $\E[\cdot\mid w_t]$:
\[
\E\!\bigl[\langle \nabla f(w_t),g_t\rangle\mid w_t\bigr]
   =\norm{\nabla f(w_t)}^{2},\qquad
\E\!\bigl[\langle \nabla f(w_t),e_t\rangle\mid w_t\bigr]=0,
\]
because $\E[e_t\mid w_t]=\E[\topk{g_t}-g_t\mid w_t]=0$ by the unbiasedness of the top‑$K$ operator in expectation (the operator is deterministic given $g_t$, hence the expectation is zero).  
Thus
\[
\E\!\bigl[\langle \nabla f(w_t),\tilde g_t\rangle\mid w_t\bigr]
   =\norm{\nabla f(w_t)}^{2}. \tag{2}
\]
[Step 2]

\noindent\textbf{Step 3.} Bound the second‑moment term $\norm{\tilde g_t}^{2}$.
\[
\begin{aligned}
\norm{\tilde g_t}^{2}
&= \norm{g_t+e_t}^{2}
 = \norm{g_t}^{2}+2\langle g_t,e_t\rangle+\norm{e_t}^{2}.
\end{aligned}
\]
Taking conditional expectation and using $\E[e_t\mid w_t]=0$,
\[
\E\!\bigl[\norm{\tilde g_t}^{2}\mid w_t\bigr]
   =\E\!\bigl[\norm{g_t}^{2}\mid w_t\bigr]
    +\E\!\bigl[\norm{e_t}^{2}\mid w_t\bigr]. \tag{3}
\]
[Step 3]

\noindent\textbf{Step 4.} Decompose $\norm{g_t}^{2}$ into variance and bias.
\[
\begin{aligned}
\E\!\bigl[\norm{g_t}^{2}\mid w_t\bigr]
&= \norm{\nabla f(w_t)}^{2}
   +\E\!\bigl[\norm{g_t-\nabla f(w_t)}^{2}\mid w_t\bigr]\\
&\le \norm{\nabla f(w_t)}^{2}+\sigma^{2}. \tag{4}
\end{aligned}
\]
[Step 4]

\noindent\textbf{Step 5.} Apply Lemma~\ref{lem:topk} to bound $\E[\norm{e_t}^{2}\mid w_t]$.
\[
\E\!\bigl[\norm{e_t}^{2}\mid w_t\bigr]
   =\E\!\bigl[\norm{\topk{g_t}-g_t}^{2}\mid w_t\bigr]
   \le\Bigl(1-\frac{K}{d}\Bigr)\E\!\bigl[\norm{g_t}^{2}\mid w_t\bigr].
\tag{5}
\]
[Step 5]

\noindent\textbf{Step 6.} Combine (3)–(5):
\[
\begin{aligned}
\E\!\bigl[\norm{\tilde g_t}^{2}\mid w_t\bigr]
&\le \E\!\bigl[\norm{g_t}^{2}\mid w_t\bigr]
   +\Bigl(1-\frac{K}{d}\Bigr)\E\!\bigl[\norm{g_t}^{2}\mid w_t\bigr] \\
&= \frac{2d-K}{d}\,\E\!\bigl[\norm{g_t}^{2}\mid w_t\bigr] \\
&\le \frac{2d-K}{d}\bigl(\norm{\nabla f(w_t)}^{2}+\sigma^{2}\bigr). \tag{6}
\end{aligned}
\]
[Step 6]

\noindent\textbf{Step 7.} Insert (2) and (6) into the smoothness inequality (1) and take total expectation.
\[
\begin{aligned}
\E\!\bigl[f(w_{t+1})\bigr]
&\le \E\!\bigl[f(w_t)\bigr]
   -\eta\E\!\bigl[\norm{\nabla f(w_t)}^{2}\bigr]
   +\frac{L\eta^{2}}{2}
      \frac{2d-K}{d}\Bigl(\E\!\bigl[\norm{\nabla f(w_t)}^{2}\bigr]
      +\sigma^{2}\Bigr). \tag{7}
\end{aligned}
\]
[Step 7]

\noindent\textbf{Step 8.} Rearrange (7):
\[
\Bigl(\eta-\frac{L\eta^{2}}{2}\frac{2d-K}{d}\Bigr)
   \E\!\bigl[\norm{\nabla f(w_t)}^{2}\bigr]
\le \E\!\bigl[f(w_t)-f(w_{t+1})\bigr]
   +\frac{L\eta^{2}}{2}\frac{2d-K}{d}\sigma^{2}. \tag{8}
\]
[Step 8]

\noindent\textbf{Step 9.} Choose $\eta$ satisfying Assumption~\ref{as:stepsize}, i.e.
\[
0<\eta\le \frac{1}{L\sqrt{1-\frac{K}{d}}}.
\]
A simple bound that holds under this choice is
\[
\eta-\frac{L\eta^{2}}{2}\frac{2d-K}{d}
\;\ge\;
\frac{\eta}{2}. \tag{9}
\]
(Proof: the left‑hand side equals $\eta\bigl(1-\frac{L\eta}{2}\frac{2d-K}{d}\bigr)$; the factor in parentheses is at least $1/2$ because $L\eta\le 1/\sqrt{1-K/d}\le 2d/(2d-K)$.)  
Insert (9) into (8):
\[
\frac{\eta}{2}\,\E\!\bigl[\norm{\nabla f(w_t)}^{2}\bigr]
\le \E\!\bigl[f(w_t)-f(w_{t+1})\bigr]
   +\frac{L\eta^{2}}{2}\frac{2d-K}{d}\sigma^{2}. \tag{10}
\]
[Step 9]

\noindent\textbf{Step 10.} Sum (10) over $t=0,\dots,T-1$ and telescope the first term.
\[
\frac{\eta}{2}\sum_{t=0}^{T-1}\E\!\bigl[\norm{\nabla f(w_t)}^{2}\bigr]
\le f(w_0)-\E\!\bigl[f(w_T)\bigr]
   +\frac{L\eta^{2}}{2}\frac{2d-K}{d}\sigma^{2}T.
\]
Since $f(w_T)\ge f^\star$, we obtain
\[
\frac{1}{T}\sum_{t=0}^{T-1}\E\!\bigl[\norm{\nabla f(w_t)}^{2}\bigr]
\le \frac{2\bigl(f(w_0)-f^\star\bigr)}{\eta T}
   + L\eta\,\frac{2d-K}{d}\sigma^{2}. \tag{11}
\]
[Step 10]

\noindent\textbf{Step 11.} Replace $\frac{2d-K}{d}=2-\frac{K}{d}$ and split the last term:
\[
L\eta\Bigl(2-\frac{K}{d}\Bigr)\sigma^{2}
=2L\eta\sigma^{2}
   -\frac{L\eta K}{d}\sigma^{2}
\le 2L\eta\sigma^{2}
   +\frac{L\eta\sigma^{2}}{K/d},
\]
where the second inequality uses $K/d\le 1$.  Hence (11) implies the bound stated in Theorem~\ref{thm:main}.

\noindent\textbf{Step 12.} Choose $\eta = \frac{c}{\sqrt{T}}$ with $c\le \frac{1}{L\sqrt{1-K/d}}$; then the first term in (11) scales as $O(1/\sqrt{T})$ and the second term as $O(1/\sqrt{T})$, establishing the $O(1/\sqrt{T})$ rate. \hfill\qedsymbol

\section*{Rate Derivation (With Constants)}
From (11) and the stepsize choice $\eta=c/\sqrt{T}$,
\[
\frac{1}{T}\sum_{t=0}^{T-1}\E\!\bigl[\norm{\nabla f(w_t)}^{2}\bigr]
\le
\frac{2\bigl(f(w_0)-f^\star\bigr)}{c}\,\frac{1}{\sqrt{T}}
\;+\;
L c\Bigl(2-\frac{K}{d}\Bigr)\sigma^{2}\,\frac{1}{\sqrt{T}}.
\]
Thus the constant multiplying $1/\sqrt{T}$ is
\[
C = \frac{2\bigl(f(w_0)-f^\star\bigr)}{c}
     + Lc\Bigl(2-\frac{K}{d}\Bigr)\sigma^{2}.
\]

\section*{Special Cases}
\begin{enumerate}
\item \textbf{Full gradient ($K=d$).}  Then $\topk{g}=g$, Lemma~\ref{lem:topk} gives zero compression error, and the bound reduces to the classic SGD result
\[
\frac{1}{T}\sum_{t=0}^{T-1}\E\!\bigl[\norm{\nabla f(w_t)}^{2}\bigr]
\le\frac{2\bigl(f(w_0)-f^\star\bigr)}{\eta T}+L\eta\sigma^{2}.
\]

\item \textbf{Extreme sparsification ($K=1$).}  The factor $1-\frac{K}{d}=1-\frac{1}{d}$ is close to $1$, so the compression term dominates; nevertheless the $O(1/\sqrt{T})$ decay still holds, albeit with a larger constant.

\item \textbf{Deterministic gradients ($\sigma=0$).}  The variance terms vanish and we obtain a deterministic rate
\[
\frac{1}{T}\sum_{t=0}^{T-1}\norm{\nabla f(w_t)}^{2}
\le\frac{2\bigl(f(w_0)-f^\star\bigr)}{\eta T},
\]
identical to gradient descent with a possibly larger stepsize bound due to compression.
\end{enumerate}

\begin{thebibliography}{9}
\bibitem{nesterov2004introductory}
Y.~Nesterov.
\newblock \emph{Introductory Lectures on Convex Optimization: A Basic Course}.
\newblock Springer Science \& Business Media, 2004.
\end{thebibliography}

\end{document}